\documentclass{article}
\usepackage[left=2cm, right=2cm, top=2cm, bottom=2cm]{geometry}
\usepackage{graphicx}
\usepackage{amsmath}
\usepackage{amssymb}
\usepackage{amsthm}
\usepackage{fancyhdr}
\usepackage{verbatim}
\usepackage{listings}
\usepackage{xcolor}
\usepackage{pdfpages}
\usepackage{cancel}
\usepackage{physics}

\lstset{
    language=C++,
    basicstyle=\ttfamily\footnotesize,
    keywordstyle=\color{blue},
    commentstyle=\color{green},
    stringstyle=\color{red},
    numbers=left,
    numberstyle=\tiny\color{gray},
    frame=single,
    breaklines=true,
    captionpos=b,
}


\title{PHYS 487: Lecture \# 14}
\author{Cliff Sun}

\newtheorem{theorem}{Theorem}[section]
\newtheorem{lemma}[theorem]{Lemma}
\newtheorem{definition}[theorem]{Definition}
\newtheorem{conjecture}[theorem]{Conjecture}
\newtheorem{proposition}[theorem]{Proposition}
\newtheorem{corollary}[theorem]{Corollary}
\newtheorem{one minute paper}[theorem]{One Minute Paper}

\pagestyle{fancy}
\lhead{\textbf{\thepage}\ \ \nouppercase{\rightmark}}
\chead{PHYS 487: Lecture \# 14}
\rhead{Cliff Sun}

\begin{document}

\maketitle

\section*{Lecture Span}
\begin{enumerate}
    \item Time dependent perturbation theory
\end{enumerate}

Dynamics: so far our Hamiltonian has been time independent which produces a phase accumulation. For example, 

\begin{equation}
    H = \hbar \omega_0 \cdot \hat{n} \cdot \frac{\vec{\sigma}}{2}
\end{equation}

Then, the wavefunction precesses around $\ket{0}$ for $\hat{n} = \hat{z}$. Consider a two-level system with 

$$
\ket{\psi(t=0)} = \ket{0}, \quad \text{initial condition}
$$

With the following Hamiltonians: 

$$H_1 = \hbar \omega_1 \frac{\sigma_z}{2}, \quad H_2 = \hbar \omega_2 \frac{\sigma_x}{2}$$ 

Now, say that we turn on the first Hamiltonian, then slowly turn the second Hamiltonian again. Then turn the first one again. In essence, this 
is simply just walking on the sphere. However, the Hamiltonians do not commute, therefore, the path of which the wavefunction takes is not the same. Recall that 

\begin{equation}
    u = e^{-\frac{i}{\hbar}\hat{H}t}
\end{equation}

Then, consider 

\begin{equation}
    u_1 = e^{-\frac{i}{\hbar}H_1 t}, \quad u_2 = e^{-\frac{i}{\hbar}H_2 t}
\end{equation}

Note that these Hamiltonians do not commute. That is:

\begin{equation}
    u_2u_1 \ket{0} \neq u_1u_2\ket{0}
\end{equation}

For now, consider a TLS (two-level system). Then, 

\begin{equation}
    H = H^{(0)} + H'(t)
\end{equation}

Here, $H^{(0)}$ is the time independent Hamiltonian and $H'(t)$ is the time dependent part. Note that $H'(t)$ is supposed to be small. Then, consider the stationary states:

\begin{equation}
    \ket{\psi_a}, \quad \ket{\psi_b}, \quad \langle \psi_i|\psi_j \rangle = \delta_{ij}, \quad H^{(0)}\ket{\psi_i} = E_i\ket{\psi_i}
\end{equation}

Moreover, there is an initial condition 

\begin{equation}
    \ket{\psi(0)} = c_a\ket{\psi_a} + c_b\ket{\psi_b}, \quad |c_a|^2 + |c_b|^2 = 1
\end{equation}

For $H' \equiv 0$, then $\ket{\psi(t)} = u \ket{\psi(0)}$ where $u$ is applied to $\ket{\psi_a}$ and $\ket{\psi_b}$. There are no "state transitions", that is 

\begin{equation}
    | \langle \psi_i | \psi(t) \rangle |^2 = |c_i|^2 = \text{const.}
\end{equation}

Next, consider the following: at $t=0$, we turn on the time dependent perturbation. Then consider the following ansatz:

\begin{equation}
    \ket{\psi(t)} = c_a(t)\ket{\psi_a}e^{-iE_a t/\hbar} + c_b(t)\ket{\psi_b}e^{-iE_bt/\hbar}
\end{equation}

Then, solve the time dependent Schrodinger wave equation: 

\begin{equation}
    H\psi(t) = i\hbar \dot{\psi(t)}, \quad \text{for $c_i(t)$} 
\end{equation}

Then insert the Ansatz into SE:

\begin{multline}
    c_a (H^{(0)} + H')\ket{\psi_a}e^{-iE_at/\hbar} + c_b(H^{(0)} + H')\ket{\psi_b}e^{-iE_bt/\hbar} \\
    = i\hbar \left[ (\dot{c_a} - c_a iE_a/\hbar)\ket{\psi_a}e^{-iE_at/\hbar} + (\dot{c_b} - c_biE_b/\hbar)\ket{\psi_b}e^{-iE_bt/\hbar} \right] 
\end{multline}

Some terms drop out. Then, multiply the entire equation by $\bra{\psi_a}$, then we obtain 

\begin{equation}
    c_a \langle \psi_a|H'|\psi_a \rangle e^{-iE_at/\hbar} + c_b\langle \psi_a | H' | \psi_b \rangle e^{-iE_bt/\hbar} = i\hbar \dot{c_a}e^{-iE_at/\hbar}
\end{equation}

Then, we obtain that 

\begin{equation}
    \dot{c_a} = -\frac{i}{\hbar}\left( c_a H'_{aa} + c_b H'_{ab}e^{i(E_a - E_b)t/\hbar}\right)
\end{equation}

Recall that $H_{ij} = \langle \psi_i | H | \psi_j \rangle$ and $A^{*}_{ij} = A_{ji}$. We obtain a similar equation for $c_b$. For now, we assume that $H_{ii}' = 0$, since the stationary states are not changed. Then, 

\begin{equation}
    \dot{c_a} = -\frac{i}{\hbar}H'_{ab}e^{-i\omega_0 t}c_b, \quad \dot{c_b} = -\frac{i}{\hbar}H'_{ba}e^{i\omega_0 t}c_a
\end{equation}

With $\omega_0 = (E_b - E_a)/\hbar$. Then, proceed with perturbation theory. Pick I.C. $c_a = 1$ and $c_b = 0$, and try to derive state transitions (none constant inner products). Then up to zero-th order, then 
$H' \rightarrow 0$ and $c_a^{(0)} = 1$ and $c_b^{(0)} = 0$. Then, we insert the 0-th order solution into the RHS of equation 14 and try to derive the first order correction of $c_a$ and $c_b$. In general, this is a \underline{recursive process}. Then 

\begin{equation}
    \dot{c_a}^{(1)} = 0 \implies c_a^{(1)}(t) =1
\end{equation}

Then with $c_b$, we have that 

\begin{equation}
    \dot{c_b}^{(1)} = -\frac{i}{\hbar}H_{ba}e^{-i\omega_0 t} \implies c_b^{(1)} = -\frac{i}{\hbar} \int_{0}^{t}dt'H'_{ba}(t')e^{i\omega_0 t'}
\end{equation}

Then, for second order, we obtain 

\begin{align}
    \dot{c_a}^{(2)} &= -\frac{i}{\hbar}H'_{ab}e^{-i\omega_0 t}\left( -\frac{i}{\hbar} \right)\int_{0}^{t}dt' \dots \\
    &= 1 - \frac{1}{\hbar^2}\int_{0}^{t}dt' H'_{ab}(t')e^{-i\omega_0 t'}\int_{0}^{t'}dt'' H'_{ba}(t'')e^{i\omega_0 t''}
\end{align}

And 

\begin{equation}
    c_b^{(2)} = c_b^{(1)}, \quad \text{because $c_a^{(1)} =1$}
\end{equation}

Sinusoidal perturbation: 

\begin{equation}
    \ket{\psi_a} = \ket{0}, \quad \ket{\psi_b} = \ket{1}, \quad H = \ket{\cdot}\bra{0} + E_1 \ket{1}\bra{1}
\end{equation}

Then consider the perturbation $H' = V \cos(\omega t)$. Then pick 

\begin{equation}
    \hat{V} = V(\ket{0}\bra{1} + \ket{1}\bra{0}) = V\sigma_x
\end{equation}

Physically, this is saying a static magnetic field along the $z$ axis and apply a sinusoidally varying b field in the x direction. Up to first order PT, we have that $c_0 =1$ and $c_1 =0$ at $t=0$. Then 

\begin{align*}
    c_1(t) &\approx -\frac{i}{\hbar}V\int_{0}^{t}\cos(\omega t')e^{i\omega_0 t'}dt' \\
    &= -\frac{iV}{2\hbar}\int_{0}^{t}dt'\left( e^{i(\omega_0 + \omega)t'} + e^{i(\omega_0 - \omega)t'} \right) \\
    &= -\frac{V}{2\hbar}\left[ \frac{e^{i(\omega_0 + \omega)t} - 1}{\omega_0 + \omega} + \frac{e^{i(\omega_0 - \omega)t} - 1}{\omega_0 - \omega} \right]
\end{align*}

This is a good approximation if $\omega \sim \omega_0$ (sorta same). Then $|\omega_0 - \omega| \ll \omega + \omega_0$. Then 

\begin{align*}
    c_1(t) &\approx -\frac{V}{2\hbar}\frac{e^{i(\omega_0 - \omega)t/2}}{\omega_0 - \omega}\left[ e^{i(\omega_0 - \omega)t/2} - e^{-i(\omega_0 - \omega)t/2}  \right] \\
    &= -i\frac{V}{\hbar}\frac{\sin(\Delta \omega t/2)}{\Delta \omega}e^{i\Delta \omega t/2}
\end{align*}

Here, $\Delta \omega = \omega_0 - \omega$. Then, the probability to find $1$ in a measurement: 

\begin{equation}
    P_{10}(t) = |c_1|^2  = \frac{V^2}{\hbar^2}\frac{\sin^2(\Delta \omega t/2)}{\Delta\omega^2}
\end{equation}



\end{document}